\chapter{配置片段库}
\chapterlead{可整段粘进 \texttt{configuration.nix} 或 \texttt{home.nix} 的常用配置。每段都注明它依赖的章节，遇到不确定的选项先用 \texttt{nixos-option} 核对。}

\section{基础}

\begin{nixcode}
{ config, pkgs, lib, ... }:

{
  # 主机与时间
  networking.hostName = "nixos-lab";
  networking.networkmanager.enable = true;
  time.timeZone = "Asia/Shanghai";
  i18n.defaultLocale = "zh_CN.UTF-8";

  # 中文显示与输入（第十六章）
  fonts.packages = with pkgs; [ noto-fonts-cjk-sans noto-fonts-cjk-serif sarasa-gothic ];
  i18n.inputMethod = {
    enable = true;
    type = "fcitx5";
    fcitx5.addons = with pkgs; [ fcitx5-chinese-addons ];
  };

  # 常用工具（第十三章）
  environment.systemPackages = with pkgs; [
    git vim curl htop jq ripgrep fd
  ];

  # 自动清理（第十五章）
  nix.gc = {
    automatic = true;
    dates = "weekly";
    options = "--delete-older-than 30d";
  };

  system.stateVersion = "26.11";
}
\end{nixcode}

\section{服务器}

\begin{nixcode}
{
  # SSH：只允许密钥登录（第二十五章）
  services.openssh = {
    enable = true;
    openFirewall = true;
    settings = {
      PasswordAuthentication = false;
      KbdInteractiveAuthentication = false;
      PermitRootLogin = "no";
      MaxAuthTries = 3;
    };
  };

  users.users.deploy = {
    isNormalUser = true;
    openssh.authorizedKeys.keys = [ "ssh-ed25519 AAAA... ci" ];
  };

  # 防火墙（第二十五章）
  networking.firewall = {
    enable = true;
    allowedTCPPorts = [ 22 80 443 ];
  };

  # 自动更新但不自动重启（第三十章）
  system.autoUpgrade = {
    enable = true;
    allowReboot = false;
    dates = "04:00";
  };

  # 防护与日志（第二十八、三十章）
  services.fail2ban.enable = true;
  services.journald.extraConfig = ''
    SystemMaxUse=1G
    MaxRetentionSec=30day
  '';
}
\end{nixcode}

\section{服务与证书}

\begin{nixcode}
{
  # 静态站点 + 自动证书（第二十六章）
  services.nginx.enable = true;
  services.nginx.virtualHosts."example.com" = {
    forceSSL = true;
    enableACME = true;
    root = "/srv/www/example";
  };
  security.acme = {
    acceptTerms = true;
    defaults.email = "you@example.com";
  };

  # 后端服务（第八章）
  systemd.services.api = {
    description = "API service";
    wantedBy = [ "multi-user.target" ];
    after = [ "network.target" ];
    serviceConfig = {
      ExecStart = "/srv/api/bin/api";
      User = "deploy";
      WorkingDirectory = "/srv/api";
      Restart = "on-failure";
      RestartSec = 3;
      StateDirectory = "api";
      MemoryMax = "1G";
    };
  };

  # 定时备份（第二十七章）
  services.restic.backups.daily = {
    repository = "/mnt/backup/restic";
    passwordFile = "/run/secrets/restic-pass";
    paths = [ "/home" "/var/lib/api" ];
    exclude = [ "/home/*/.cache" ];
    timerConfig = { OnCalendar = "daily"; Persistent = true; };
    pruneOpts = [ "--keep-daily 7" "--keep-weekly 4" "--keep-monthly 6" ];
  };
}
\end{nixcode}

\section{桌面}

\begin{nixcode}
{
  # Hyprland（第十六章）
  programs.hyprland.enable = true;
  services.displayManager.ly.enable = true;

  # 状态栏、通知、截图
  programs.waybar.enable = true;
  services.mako.enable = true;
  environment.systemPackages = with pkgs; [ grim slurp wl-clipboard ];

  # 硬件（第十七章）
  hardware.bluetooth.enable = true;
  hardware.enableRedistributableFirmware = true;
  services.tlp.enable = true;
  services.fwupd.enable = true;
}
\end{nixcode}

\section{Home Manager}

\begin{nixcode}
{ pkgs, ... }:

{
  home.username = "weiwen";
  home.homeDirectory = "/home/weiwen";
  home.stateVersion = "26.11";

  programs.bash.enable = true;
  programs.git = {
    enable = true;
    userName = "Weiwen Zhang";
    userEmail = "you@example.com";
    extraConfig = {
      init.defaultBranch = "main";
      pull.rebase = true;
      core.pager = "delta";
    };
  };
  programs.zoxide = { enable = true; enableBashIntegration = true; };
  programs.starship.enable = true;

  home.packages = with pkgs; [ ripgrep fd bat eza jq ];
}
\end{nixcode}
